\documentclass[preview]{standalone}

\usepackage{amsmath}
\usepackage{amssymb}
\usepackage{stellar}
\usepackage{definitions}

\begin{document}

\id{minimal-polynomial-canonical-forms-exercises}
\genpage

\section{Exercises}

\begin{snippetexercise}{canonical-form-over-rationals-and-complex-exercise}{Canonical form over the rationals and the complexes}
    Determine the canonical form of
    \[
        A=
        \begin{pmatrix}
            1&0&0&0\\
            2&1&0&0\\
            0&0&-\frac12&-\frac32\\
            0&0&\frac12&-\frac12
        \end{pmatrix}
    \]
    first over \(\rationalnumbers\), and then over
    \(\complexnumbers\).
\end{snippetexercise}

\begin{snippetsolution}{canonical-form-over-rationals-and-complex-solution}{Canonical form over the rationals and the complexes}
    The characteristic polynomial is
    \[
        \chi_A(X)=(X-1)^2(X^2+X+1).
    \]
    The second factor is irreducible in \(\rationalnumbers[X]\). Over
    \(\rationalnumbers\), the primary decomposition is
    \[
        \rationalnumbers[X]/((X-1)^2)
        \oplus
        \rationalnumbers[X]/(X^2+X+1).
    \]
    Since the two annihilators are coprime, the associated module is
    cyclic:
    \[
        V_A
        \moduleisomorphic
        \rationalnumbers[X]/\bigl((X-1)^2(X^2+X+1)\bigr).
    \]
    Hence the rational canonical form is the companion matrix
    \[
        C\bigl((X-1)^2(X^2+X+1)\bigr).
    \]

    Over \(\complexnumbers\),
    \[
        X^2+X+1=(X-\omega)(X-\omega^2),
        \qquad
        \omega=-\frac12+\frac{\sqrt3}{2}i.
    \]
    Since \(\mrank(A-I_4)=3\), the eigenspace for \(1\) has dimension
    \(1\), so its algebraic multiplicity \(2\) produces one block of
    size \(2\). The other eigenvalues are simple. Therefore
    \[
        J(A)=J_2(1)\oplus J_1(\omega)\oplus J_1(\omega^2).
    \]
\end{snippetsolution}

\begin{snippetexercise}{minimal-polynomial-jordan-form-parameter-exercise}{Minimal polynomial and Jordan form with a parameter}
    For \(t\in\rationalnumbers\), let
    \[
        A_t=
        \begin{pmatrix}
            1&0&1&0\\
            0&1&0&0\\
            0&0&5&3\\
            0&t&-3&-1
        \end{pmatrix}.
    \]
    Determine the minimal polynomial and the Jordan form of \(A_t\) as
    \(t\) varies.
\end{snippetexercise}

\begin{snippetsolution}{minimal-polynomial-jordan-form-parameter-solution}{Minimal polynomial and Jordan form with a parameter}
    For every \(t\),
    \[
        \chi_{A_t}(X)=(X-1)^2(X-2)^2.
    \]
    At the eigenvalue \(1\),
    \[
        \mrank(A_t-I_4)
        =
        \begin{cases}
            2,&t=0,\\
            3,&t\neq0.
        \end{cases}
    \]
    Thus \(t=0\) gives two \(1\times1\) blocks for the eigenvalue \(1\),
    while \(t\neq0\) gives one block \(J_2(1)\).

    At the eigenvalue \(2\), one has
    \[
        \mrank(A_t-2I_4)=3
    \]
    for every \(t\), so the \(2\)-primary part is always \(J_2(2)\).
    Consequently,
    \[
        J(A_t)=
        \begin{cases}
            J_1(1)\oplus J_1(1)\oplus J_2(2),&t=0,\\
            J_2(1)\oplus J_2(2),&t\neq0,
        \end{cases}
    \]
    and
    \[
        \mu_{A_t}=
        \begin{cases}
            (X-1)(X-2)^2,&t=0,\\
            (X-1)^2(X-2)^2,&t\neq0.
        \end{cases}
    \]
\end{snippetsolution}

\begin{snippetexercise}{canonical-form-over-rationals-and-reals-exercise}{Canonical form over the rationals and the reals}
    Let
    \[
        A=
        \begin{pmatrix}
            1&0&0&-2&-1\\
            1&0&0&2&1\\
            0&1&2&1&0\\
            -1&2&0&0&-1\\
            1&-2&0&2&3
        \end{pmatrix}.
    \]
    Determine the elementary divisors and canonical form of \(A\) over
    \(\rationalnumbers\) and over \(\realnumbers\).
\end{snippetexercise}

\begin{snippetsolution}{canonical-form-over-rationals-and-reals-solution}{Canonical form over the rationals and the reals}
    The characteristic polynomial is
    \[
        \chi_A(X)=(X-2)^3(X^2-2).
    \]
    Moreover,
    \[
        \mrank(A-2I_5)=3,
        \qquad
        \dim\ker(A-2I_5)=2.
    \]
    Thus the eigenvalue \(2\), of algebraic multiplicity \(3\), has two
    blocks of sizes \(2\) and \(1\).

    Over \(\realnumbers\), the other two eigenvalues are the simple
    roots \(\sqrt2\) and \(-\sqrt2\). Hence
    \[
        J(A)
        =
        J_1(\sqrt2)\oplus J_1(-\sqrt2)
        \oplus J_2(2)\oplus J_1(2).
    \]

    Over \(\rationalnumbers\), the factor \(X^2-2\) is irreducible. The
    elementary-divisor decomposition is
    \[
        V_A
        \moduleisomorphic
        \rationalnumbers[X]/(X^2-2)
        \oplus
        \rationalnumbers[X]/((X-2)^2)
        \oplus
        \rationalnumbers[X]/(X-2).
    \]
    Right-aligning the primary cyclic factors gives the invariant
    factors
    \[
        d_1=X-2,
        \qquad
        d_2=(X^2-2)(X-2)^2.
    \]
    Therefore
    \[
        \mu_A=(X^2-2)(X-2)^2
    \]
    and the rational canonical form is
    \[
        C(X-2)\oplus C\bigl((X^2-2)(X-2)^2\bigr).
    \]
    The second companion block corresponds to
    \[
        (X^2-2)(X-2)^2
        =
        X^4-4X^3+2X^2+8X-8.
    \]
\end{snippetsolution}

\end{document}
